%! Author = charon
%! Date = 6/4/24
\section{Skills}\label{sec:kenntnisse}
% This section is laid out using a table. A \tableentry command adds lines with the following parameters:
%\tableentry{Heading}{Content}{spaceafter}
% All 3 parameters must be supplied but any can be empty if you don't need them
% A "spaceafter" value in the third parameter will add some vertical space -- this is to be used between headings, leave it empty for no extra space
%------------------------------------------------
\begin{supertabular}{r l} % Start a table with two columns, the table will ensure everything is aligned
    \tableentry{Sprachen}{Deutsch (Muttersprache),}{spaceafter}
    \tableentry{}{Tschechisch (Muttersprache),}{spaceafter}
    \tableentry{}{Englisch (C1+), Französisch (B1)}{spaceafter}
    \tableentry{Technische Skills}{\textit{Programmierung}}{spaceafter}
    \tableentry{}{Python, Java, Kotlin, Javascript,}{spaceafter}
    \tableentry{}{Swift, Flutter, C, Rust}{spaceafter}
    \tableentry{}{\textit{Security}}{spaceafter}
    \tableentry{}{Netzwekanalyse, Fuzzing,} {spaceafter}
    \tableentry{}{CTFs, Forschung}{spaceafter}
    \tableentry{}{\textit{Netzwerke}}{spaceafter}
    \tableentry{}{CCNA, TCP/IP Analyse}{spaceafter}
    \tableentry{}{\textit{Systeme}}{spaceafter}
    \tableentry{}{Linux (Administration \& Hardening),}{spaceafter}
    \tableentry{}{Proxmox, Docker, Kubernetes,}{spaceafter}
    \tableentry{}{Git, CI/CD}{spaceafter}
\end{supertabular}
